% ============================================================
% CAP 4 — ANÁLISIS DE MERCADO Y TENDENCIAS
% ============================================================

\section{An\'alisis de mercado y tendencias}
\label{sec:mercado-tendencias}

\subsection{Benchmarking de servicios de movilidad nocturna en corredores análogos}
\label{sec:benchmarking}

Esta sección analiza casos documentados de servicios
de movilidad nocturna en corredores con demanda
concentrada y predecible, comparables al corredor
CHUF--Ferrol en escala o perfil de usuario. El objetivo
no es seleccionar un modelo: es identificar qué
condiciones de operación, financiación y aceptación
han producido resultados positivos y qué factores han
limitado la viabilidad. La evaluación de candidatos
frente a los criterios del presente TFM se realiza
en el capítulo~\ref{cap:propuesta}.

\subsubsection{Viabilidad económica del bus fijo nocturno: umbral de demanda}
\label{sec:umbral-demanda}

El punto de partida del benchmarking es una cuestión
económica fundamental: ¿a partir de qué nivel de demanda
resulta viable operar una línea de transporte público
fija en franja nocturna? \textcite{Estrada2021} analizan
mediante un modelo analítico calibrado con datos reales
del área metropolitana de Barcelona tres modalidades de
servicio (bus de ruta fija con paradas variables,
microbús de ruta flexible y taxi compartido) en función
de la densidad de demanda. Sus resultados indican que
el bus de ruta fija minimiza el coste total por pasajero
únicamente a partir de densidades superiores a
200~pasajeros/km²·h, mientras que la demanda nocturna
típica en entornos urbanos se sitúa en torno a
2~pasajeros/km²·h \parencite{Estrada2021}. Esta brecha
de dos órdenes de magnitud explica estructuralmente
por qué los operadores retiran el servicio convencional
en horario nocturno y abre el espacio para modelos
de servicio alternativos.

\subsubsection{DRT nocturno en ciudad media: el caso de Padua}
\label{sec:caso-padua}

El servicio QuiBus de Padua (Italia, aproximadamente
210.000~habitantes) constituye uno de los casos europeos
mejor documentados de transporte a demanda nocturno en
ciudad de tamaño medio. Operado por Busitalia Veneto y
financiado conjuntamente por el Ayuntamiento de Padua
y la Universidad de Padua, el servicio cubre la franja
21:00--24:00 de domingo a jueves y 21:00--03:00 los
viernes y sábados. Los usuarios reservan el trayecto
mediante aplicación móvil indicando parada de origen,
hora de llegada deseada y número de pasajeros
\parencite{Busitalia2019QuiBus}.

Los resultados operativos del piloto son relevantes:
7.000~usuarios registrados en el primer mes y 19.000
en los primeros seis meses, con un ritmo de crecimiento
de 300~nuevos registros por semana y una valoración
media de 4{,}7 sobre~5 en la aplicación. El crecimiento
sostenido llevó a prorrogar el servicio más allá del
período piloto inicial \parencite{Busitalia2019QuiBus}.
El modelo de financiación público-institucional y la
sincronización con la demanda universitaria nocturna
son los dos factores que los propios operadores
identifican como determinantes del éxito.

\subsubsection{Shuttle DRT para personal sanitario: el caso de Hull}
\label{sec:caso-hull}

El caso de Hull (Reino Unido) documenta un modelo de
shuttle a demanda diseñado específicamente para personal
sanitario de turno. Durante la reducción de la red
regular de autobuses en 2020, Stagecoach Solutions y
el Hull University Teaching Hospitals NHS Trust
implantaron un servicio de lanzaderas con reserva
mediante aplicación, operando en las franjas de
madrugada y noche para acomodar todos los patrones
de turno hospitalario \parencite{Stagecoach2020Hull}.

El servicio utilizó vehículos de piso bajo de tamaño
estándar, con un trayecto medio de aproximadamente
20~minutos y 8{,}8~km. La financiación se articuló
mediante una ronda inicial de fondos públicos del
Ayuntamiento de Hull sobre la que se construyó el
esquema operativo continuo. La colaboración directa
entre operador de transporte e institución sanitaria
como promotora del servicio es el elemento diferencial
de este caso respecto a los DRT urbanos genéricos
\parencite{Stagecoach2020Hull}.

\subsubsection{DRT nocturno y trabajadores de turno: evidencia del mapa europeo}
\label{sec:drt-mapa-europeo}

El informe de mapeo de servicios DRT en Europa del
K2 (Swedish Knowledge Centre for Public Transport)
documenta que el servicio MK~Connect de Milton Keynes
(Reino Unido) opera con horarios extendidos hasta las
02:00 en determinadas zonas, y es valorado de forma
explícita por los trabajadores de turno en hospitales
locales como alternativa al vehículo privado en
franjas sin cobertura convencional. La financiación
está integrada en el sistema tarifario del transporte
público general \parencite{K2Centrum2025DRT}.

El informe concluye que los servicios DRT en ciudades
medias europeas muestran resultados variables en función
de tres factores: el modelo de financiación, la tasa
de adopción digital por parte de los usuarios objetivo,
y el grado de integración con los horarios de los
principales generadores de demanda nocturna
\parencite{K2Centrum2025DRT}.

\subsection{Síntesis del benchmarking}
\label{sec:sintesis-benchmarking}

Los casos analizados permiten extraer cuatro
observaciones transferibles al corredor CHUF--Ferrol:

\begin{enumerate}
  \item La demanda nocturna en ciudades medias es
    estructuralmente insuficiente para justificar
    una línea de bus fija con criterios de rentabilidad
    convencional \parencite{Estrada2021}.

  \item Los modelos DRT con reserva digital han
    demostrado viabilidad operativa en ciudades de
    tamaño comparable a Ferrol cuando existe un
    generador de demanda concentrada y predecible
    \parencite{Busitalia2019QuiBus, K2Centrum2025DRT}.

  \item La colaboración entre el operador de transporte
    y la institución con personal de turno nocturno
    (hospital, universidad) es un factor común en los
    casos con financiación estable y continuidad tras
    el piloto \parencite{Stagecoach2020Hull,
    Busitalia2019QuiBus}.

  \item La sincronización del servicio con los horarios
    de los turnos (no con frecuencias genéricas de
    red) es el criterio de diseño operativo que los
    usuarios de turno identifican como más relevante
    \parencite{K2Centrum2025DRT, Stagecoach2020Hull}.
\end{enumerate}

Estas observaciones se incorporan como criterios
de evaluación en el capítulo~\ref{cap:criterios},
donde se contrastan con los insights derivados de
la encuesta propia \parencite{LouridoRegueira2026}.


\subsection{Análisis de tendencias}
\label{sec:tendencias}
\begin{pendiente}
\end{pendiente}

\subsubsection{Tendencias políticas}
\label{sec:tendencias-politicas}
\begin{pendiente}
\end{pendiente}

\subsubsection{Tendencias tecnológicas}
\label{sec:tendencias-tecnologicas}
\begin{pendiente}
\end{pendiente}

\subsubsection{Tendencias sociales}
\label{sec:tendencias-sociales}
\begin{pendiente}
\end{pendiente}

\subsubsection{Tendencias de movilidad y consumo}
\label{sec:tendencias-movilidad}
\begin{pendiente}
\end{pendiente}
